\chapter{IMAGE VERBALIZATION}
\label{chap:2nd chap}

% ==================================================================
\section{Introduction}
\label{sec:ch2_intro}
% ==================================================================

A drone operating in an unstructured indoor environment must continuously interpret its surroundings and decide the next manoeuvre. 
This decision requires understanding not just what the camera
sees, but what it means in the context of the flight: identical depth sensor readings from a nearby wall and a nearby person may demand completely different pilot actions. 
Rule-based systems that threshold sensor values cannot make this distinction. 
Human language, on the contrary, is inherently context-sensitive. Large language models (LLMs), trained in billions of words of grounded description, possess exactly this kind of contextual reasoning capability \cite{vemprala2023chatgpt}.

The practical obstacle is the camera. A small, weight-constrained indoor drone cannot carry a high-resolution sensor; the onboard camera in this work is an ESP32-S3-Sense module streaming JPEG frames at \textbf{320$\times$240~pixels} (approximately 10~KB per frame). 
Such frames are suitable for compressed transmission over the companion computer link but are far below the image resolution at which LLM vision has been benchmarked \cite{achiam2023gpt4,
reid2024gemini, liu2023llava}.
\newpage
Three questions motivate this chapter:
\begin{enumerate}
  \item Can an LLM produce a meaningful scene description and a pilot
        action suggestion from a 320$\times$240 JPEG that costs only 85--102 input tokens for the image itself?
  \item What prompt structure, token budget, and context history minimise the risk of a dangerous recommendation?
  \item Does a layered architecture, combining a millisecond-scale local emergency detector, a sub-second structural vision layer, and the LLM cognitive layer, provide higher action safety than any single component alone?
\end{enumerate}

This chapter answers all three. 
A systematic validation series (\ref{sec:ch2_validation}) optimises the LLM call parameters such as prompt technique, output token budget, scene context history, and sampling temperature using controlled single-variable experiments on the eight canonical drone
scenes captured from real ESP32-S3-Sense camera frames.
The integration series (\ref{sec:ch2_gseries}) then assembles all components and proves that the proposed layered architecture achieves \textbf{97.5\%~action safety} across 157~trials, with two models reaching 100\%. 
A hardware validation placeholder (\ref{sec:ch2_h5}) outlines the
live-flight experiment that closes the simulation-to-hardware gap.


% ==================================================================
\section{Problem Statement}
\label{sec:ch2_problem}
% ==================================================================

\subsection{The Visual-Semantic Gap in Low-Resolution Drone Imagery}
\label{sec:ch2_problem_gap}

Indoor drone cameras face a fundamental constraint, the companion computer budget (weight, power, cost) limits the sensor to low-cost modules whose native output is a compressed JPEG at VGA or below. The ESP32-S3-Sense camera used in this work outputs \textbf{320$\times$240~pixels at approximately 9.7~KB per JPEG frame} (measured across five frames per scene;
Table~\ref{tab:ch2_image_spec}). 
This resolution is sufficient for human navigation in controlled environments but introduces three failure modes for
automated reasoning:
\newpage
\begin{enumerate}
  \item \textbf{Feature ambiguity.} At 320$\times$240, a featureless grey wall filling the frame and an empty grey doorway can be visually indistinguishable. 
  Standard object detectors report ``no objects detected'' for both; only spatial reasoning over texture gradients or depth cues can separate them.
  \item \textbf{Context loss.} A person at 3~m occupies $\sim$20~pixels of vertical extent; a person at 0.3~m occupies $\sim$170~pixels. 
  A fixed-threshold person detector trained on standard resolution images will fail to fire at one extreme or the other, \cite{persondetectionuav2025}.
  \item \textbf{Blocked-sensor blindness.} If the lens is partially covered (dust, hand, water), a standard detection pipeline interprets zero detections as ``clear path'' and commands forward motion directly into the obstruction.
\end{enumerate}

\begin{table}[ht]
  \caption{Image payload specification: ESP32-S3-Sense camera used in all
    V-series and G-series experiments. Token counts are per-frame costs for the
    vision LLM input: GPT-4o and GPT-4o-mini in low-detail mode (fixed 85
    tokens); Claude (Sonnet) by the formula $(w \times h)/750$; Gemini uses
    per-character billing rather than token counts, making direct comparison
    approximate. Values measured from five representative frames.}
  \label{tab:ch2_image_spec}
  \centering
  \small
  \begin{tabular}{ll}
    \toprule
    Property & Value \\
    \midrule
    Sensor module & ESP32-S3-Sense \\
    Resolution & 320$\times$240~pixels \\
    Format & JPEG (compressed) \\
    Mean file size & 9.7~KB (range: 7.1--12.4~KB) \\
    Mean pixel count & 76,800 \\
    \midrule
    Image token cost --- GPT-4o / GPT-4o-mini & 85~tokens (low-detail, fixed) \\
    Image token cost --- Claude (Sonnet) & $\approx$102~tokens ($(320 \times 240)/750$) \\
    Image token cost --- Gemini & $\approx$flat (pixel-area billing) \\
    \midrule
    Full prompt input tokens (with YOLO metadata) & 575--1,008 per call \\
    Image fraction of total input cost & $\sim$9--15\% \\
    \bottomrule
  \end{tabular}
\end{table}

As Table~\ref{tab:ch2_image_spec} shows, the 320$\times$240 image itself accounts for only 85--102 tokens which is less than 15\% of a typical API call.
The dominant cost is the text prompt (system instruction, scene context history, sensor metadata). 
This means that the vision payload is not the cost bottleneck; the system is open to the per-call budgeting and cost optimisation demonstrated in the Integration series experiments.

\subsection{Sensor-Only Approaches and Their Limits}
\label{sec:ch2_problem_sensors}

Existing approaches to indoor drone safety combine depth sensors, IMUs, and rule-based thresholds. 
Time-of-Flight (ToF) sensors report a scalar distance that is identical whether the obstacle is a wall, a person, or a box.
YOLO-based detectors trained on standard COCO datasets achieve $>$80\% precision on person detection at typical resolutions but degrade significantly on UAV-perspective imagery \cite{persondetectionuav2025,kim2024yoloihd}. 
Rule engines that map detection categories to actions inherit every false positive and false negative from the detector layer with no mechanism for contextual recovery. 
Critically, no rule-based system can detect that its own sensor has
failed: a hand over the lens produces zero detections, and a threshold engine concludes ``safe to advance''.

Language model-equipped drone systems that do not address the low-resolution problem either operate offline, assume high-bandwidth video, or bypass the image entirely by relying on structured state descriptions \cite{ahn2022saycan, huang2022innermonologue}. 
Recent systems that do incorporate on-board cameras (e.g., \cite{rescuedrone2025, hierarchicaldrone2025}) operate on higher
resolution feeds without the strict 10~KB frame-size constraint of a
weight-limited nano-drone. 
This chapter directly targets the low-resolution, token-budget-constrained regime.
\newpage
\subsection{Research Questions and Scope}
\label{sec:ch2_problem_rq}

Given the constraints above, the contribution of this chapter is scoped as follows:

\begin{description}
  \item[RQ1 - Feasibility:] Can an off-the-shelf multimodal LLM produce a correct pilot action recommendation from a 320$\times$240 JPEG with a tructured text prompt, even when image quality prevents unambiguous visual identification of the hazard?
  \item[RQ2 - Optimisation:] What prompt technique, output token budget, and scene context history minimise dangerous recommendations while remaining within a practical API cost envelope?
  \item[RQ3 - Architecture:] Does a layered hierarchy: local emergency detector, structural vision layer and a LLM cognitive layer outperform any single layer alone and does real-time emergency interruption improve hazard-detection timing without increasing cost?
\end{description}


% ==================================================================
\section{Image Verbalization with Layered Vision}
\label{sec:ch2_proposed}
% ==================================================================

\subsection{Core Idea: Augmenting the Image with Structured Metadata}
\label{sec:ch2_proposed_idea}

The central insight is that a low-resolution image alone may be insufficient for reliable classification, but an image paired with structured sensor metadata like object class, bounding box, depth estimate, and a proximity warning flag provides enough grounding for an LLM to reason safely. 
The image contributes global scene context (lighting, texture
fill, apparent distance, room layout); the metadata contributes local precise measurements that the LLM uses to cross-check its visual analysis.

This enrichment process is called an image verbalization in the context of this thesis. 
A camera frame is converted into a structured natural-language description by combining visual analysis (the LLM looking at the JPEG) with metadata extracted by fast local detectors (YOLO-World, DepthAnything~v2, MediaPipe EfficientDet-Lite0). 
The verbalizer then produces a risk label (\texttt{safe}, \texttt{caution}, or \texttt{hazard}) and a pilot action suggestions such as (\texttt{PITCH\_FORWARD}, \texttt{HOVER}, or \texttt{LAND/STOP}).

\subsection{Four-Tier Layered Architecture}
\label{sec:ch2_proposed_tiers}

A single LLM call cannot serve as the efficient scene description layer, at 2,333--7,218~ms per call, the LLM operates at 0.1--0.4~Hz while a drone moving at 0.5~m/s travels 12--36~cm between calls. 
There is a high probability that scheduled LLM triggers might miss to capture emergency scenes between schedules.
Emergency events require faster detection. 
The proposed architecture separates perception and control
into four levels, each operating at the highest frequency its computational cost permits (Figure~\ref{fig:ch2_concept}):

\begin{description}
  \item[Tier~1 - PID Firmware (2~kHz):]
    The flight controller's PID loop maintains attitude, altitude, and position hold. 
    It receives set-points from the higher tiers and executes at 2,000~Hz with sub-millisecond latency. 
    The LLM never queries this tier or waits for an API response.
    \newpage
  \item[Tier~1.5 - MediaPipe Emergency Detector ($\sim$60~Hz):]
    MediaPipe EfficientDet-Lite0 \cite{tan2020efficientdet, lugaresi2019mediapipe} runs on every camera frame in $\sim$16~ms with no API cost. 
    It performs three binary checks like person detection (conf $>$ 0.30), wall-fill texture check (gradient standard deviation $<$ 30), and brightness gate (mean luminance $<$ 40). 
    Any positive check fires a real-time interrupt that triggers an immediate LLM call, bypassing the scheduled polling interval.
    This tier handles emergencies between scheduled cognitive-tier updates.
  \item[Tier~2 - YOLO-World + DepthAnything~v2 ($\sim$4~Hz):]
    YOLO-World \cite{cheng2024yoloworld} provides open-vocabulary structural detection across 17 indoor hazard classes (walls, doors, furniture, persons). 
    DepthAnything~v2 Metric~Indoor \cite{yang2024depthanythingv2}
    provides absolute metric depth estimates in metres. Together they produce the sensor metadata string that accompanies every LLM call, run at $\sim$275~ms per frame. 
    A texture-based proximity flag is appended when the detected depth is inconsistent with the wall-fill visual signature, guarding against the DepthAnything~v2 known failure mode on featureless surfaces.
    Primary purpose of this tier is to pass structured metadata to the LLM.
  \item[Tier~3 --- LLM Cognitive Reasoner (0.1--0.4~Hz):]
    The LLM receives the current JPEG frame plus a structured text prompt containing the sensor metadata and the two most recent prior frame descriptions (short context history). 
    It produces a structured JSON response containing scene description, proximity, risk level, reasoning, and recommended action. 
    The LLM will describe the scene and suggest the recommended action to the human operator. 
    The human operator can accepts or override every recommendation.
\end{description}

This architecture ensures that no API response is on the critical latency path for basic attitude control (Tier~1), emergency event trigger (Tier~1.5), or structural metadata provider (Tier~2). 
The LLM contributes contextual reasoning at the cognitive timescale without blocking the faster layers.

% -------------------------------------------------------------------
\begin{figure}[ht]
  \centering
  \includegraphics[width=\textwidth]{fig2_2_verbalization_concept.pdf}
  \caption{Image verbalization architecture.}
  \label{fig:ch2_concept}
\end{figure}
% --------

\newpage
% ==================================================================
\section{Experimental Setup}
\label{sec:ch2_setup}
% ==================================================================

\subsection{Reference Image Scenes}
\label{sec:ch2_setup_scenes}

All validation series and integration series experiments use the same eight canonical scenes captured from the ESP32-S3-Sense camera in the same indoor environment (laboratory with tables, shelving, chairs, and controlled lighting).
Each scene was photographed in five independent repetitions (different times of the day, slight variation in viewpoint) to generate 40~frames per scene (320 total in all 8~scenes). 
Figure~\ref{fig:ch2_scene_grid} shows one representative
frame per scene; Table~\ref{tab:ch2_scenes} gives the ground-truth label and the rationale for this label.

% -------------------------------------------------------------------
\begin{figure}[ht]
  \centering
  \includegraphics[width=\textwidth]{fig2_1_scene_grid.pdf}
  \caption{Eight canonical reference scenes used in all validation series and integration series
    experiments. (a). Camera Lens blocked. (b). Cluttered obstacles. (c). Laptop on a table. (d). Person near. (e). Lights dimmed. (f). Opened Door. (g). Wall at 25 cm. (h). Person far.}
  \label{fig:ch2_scene_grid}
\end{figure}
% -------------------------------------------------------------------

\begin{table}[ht]
  \caption{Eight canonical scenes: ground-truth label, sensor signature, and principal challenge for the vision pipeline.
  All scenes captured at 320$\times$240 from ESP32-S3-Sense.}
  \label{tab:ch2_scenes}
  \centering
  \small
  \begin{tabular}{llll}
    \toprule
    Scene & Truth & Sensor signature & Principal challenge \\
    \midrule
    \texttt{person\_near}  & hazard  & Person conf $>$ 0.4, depth $<$ 0.5~m & Correct label; priority escalation \\
    \texttt{wall\_close}   & hazard  & DA~v2 may report 2.09~m (wrong)        & Depth sensor error; visual fill test \\
    \texttt{blocked\_lens} & hazard  & Zero YOLO detections; low brightness   & Rule-based engines return ``safe'' \\
    \texttt{person\_far}   & caution & Person conf $\approx$ 0.3, depth $>$ 2~m & Over-caution expected; must be safe \\
    \texttt{dim\_light}    & caution & Low mean luminance ($\approx$40--80)   & Ambiguous: models often say hazard \\
    \texttt{cluttered}     & caution & Many YOLO objects; no single dominant  & Models over-classify as hazard \\
    \texttt{object\_table} & caution & Table/laptop obstructing lower portion & Partial blockage at flight altitude \\
    \texttt{door\_open}    & safe    & Low object density; open passage ahead & Must not HOVER unnecessarily \\
    \bottomrule
  \end{tabular}
\end{table}

\newpage
The scene set covers the three risk categories (\texttt{hazard},
\texttt{caution}, \texttt{safe}) in a 3:4:1 ratio, with deliberate inclusion of adversarial cases: \texttt{blocked\_lens} tests sensor-failure detection; \texttt{wall\_close} tests depth-sensor disagreement with visual evidence; \texttt{person\_far} tests resistance to over-caution; \texttt{door\_open} tests resistance to unnecessary hovering.

\subsection{Hardware and Software Configuration}
\label{sec:ch2_setup_hw}

The camera frames were captured from the ESP32-S3-Sense streaming server over a local Wi-Fi link to the companion computer.
The Pre-processing applies CLAHE contrast enhancement \cite{clahe_zuiderveld1994} (\texttt{clip\_limit} = 2.0, \texttt{tile\_grid} = (8,8)) before all downstream inferences. 
YOLO-World and DepthAnything~v2~Metric~Indoor run on CPU without
hardware acceleration. 
All LLM API calls use the models' multimodal endpoints.
The experiment parameters for each run of the validation series are given in the corresponding subsection; the shared baseline between all experiments is five independent runs per scene per condition unless otherwise stated \cite{liang2023helm}.

\begin{table}[ht]
  \caption{LLM models evaluated across validation series and integration series experiments.
  Costs are measured (not listed prices) from the experiment CSV files.
  Latency reported as mean across all successful calls per model in}
  \label{tab:ch2_models}
  \centering
  \small
  \begin{tabular}{lllrr}
    \toprule
    Model & Tier & API & Mean latency & Est.\ cost/call \\
    \midrule
    Claude (Sonnet)   & 3 & Anthropic & 7,777~ms & \$0.010 \\
    GPT-4o            & 3 & OpenAI    & $\sim$2,910~ms & \$0.008 \\
    GPT-4o-mini       & 3 & OpenAI    & 4,196~ms & \$0.001 \\
    \textbf{Gemini} (2.5-Flash)  & 3 & Google  & \textbf{2,554~ms} & \textbf{\$0.0002} \\
    \bottomrule
  \end{tabular}
\end{table}


% ==================================================================
\section{Validation Experiments: LLM Verbalization Parameters}
\label{sec:ch2_validation}
% ==================================================================

Before integrating all tiers into a single pipeline, three LLM call parameters are optimised through controlled single-variable experiments. 
The validation series isolates each parameter while keeping all other parameters fixed. 
Each experiment uses the same eight canonical scenes in five runs per condition. 
The metric hierarchy is: (1) action safety (was the pilot action recommended safe given ground truth?), (2) dangerous case count (truth~= ~ hazard and action~= ~ PITCH \_FORWARD), (3) label precision (did the LLM assign the correct risk category?) and
(4) verbalization quality (0--5 rubric that assesses completeness, grounding and actionability). 
Action safety and dangerous case count are the primary safety metrics; label accuracy is secondary for LLMs.

% ------------------------------------------------------------------
\subsection{Baseline Model Comparison}
\label{sec:ch2_v1}
% ------------------------------------------------------------------

This experiment establishes baseline performance across all four models using a combined prompt that includes YOLO-World metadata, MediaPipe metadata, and CLIP scene labels. 
All models receive identical inputs (same JPEG + same metadata).
Table~\ref{tab:ch2_v1} reports per-model baseline results across 160~trials (4~models $\times$ 8~scenes $\times$ 5~runs).

\begin{table}[ht]
  \caption{V1 baseline model comparison: all models on the same 8 scenes × 5 runs = 160~trials. 
  Combined prompt (YOLO + MediaPipe + CLIP). All models achieve 100\% action safety at baseline; differences in label accuracy and quality motivate the subsequent parameter studies.}
  \label{tab:ch2_v1}
  \centering
  \small
  \begin{tabular}{lrrrrr}
    \toprule
    Model & $N$ & LblAcc & Quality/5 & Mean latency & Dangerous \\
    \midrule
    Claude       & 40 & 62.5\% & 3.9 & 9,069~ms & 0 \\
    GPT-4o       & 40 & \textbf{70.0\%} & \textbf{4.4} & 3,078~ms & 0 \\
    GPT-4o-mini  & 40 & 57.5\% & 4.0 & 4,836~ms & 0 \\
    Gemini       & 40 & 62.5\% & 4.1 & \textbf{2,468~ms} & 0 \\
    \midrule
    All combined & 160 & 63.1\% & 4.1 & 4,863~ms & \textbf{0} \\
    \bottomrule
  \end{tabular}
\end{table}

The V1 baseline confirms that all four models can reason safely about
320$\times$240 frames and produce zero dangerous recommendations under the combined prompt. 
GPT-4o achieves the highest label accuracy (70.0\%) and
verbalization quality (4.4/5); Gemini achieves the lowest latency (2,468~ms).
The fact that all models achieve 100\% action safety at baseline validates the core feasibility hypothesis (RQ1): a multimodal LLM can produce safe pilot action suggestions from a 320$\times$240 drone camera image.

% ------------------------------------------------------------------
\subsection{Prompt Technique}
\label{sec:ch2_v2}
% ------------------------------------------------------------------

Five prompt techniques are compared: zero-shot (no examples), few-shot-3 (three labelled examples), chain-of-thought at 300~tokens (CoT-300), chain-of-thought at 600~tokens (CoT-600), and a structured JSON format that enforces a fixed output schema. 
The structured format architecturally prevents the most dangerous failure mode: a model that internally assesses a scene as
hazard cannot then output \texttt{PITCH\_FORWARD}, because the schema checks \texttt{risk\_level} and \texttt{recommended\_action} as separate fields. 
A contradictory pair is immediately detectable as a parse error.

The experiment runs 790~valid trials across 4~models, 4-5~techniques, 8~scenes, 5~runs; results are shown in
Table~\ref{tab:ch2_v2_technique} and Figure~\ref{fig:ch2_v2_techniques}.

\begin{table}[ht]
  \caption{V2 prompt technique comparison: all four models combined, 790 valid
    trials. LblAcc = label accuracy; DescAcc = description accuracy; RsnRisk =
    reasoning-risk alignment; RsnAct = reason-action alignment; ActSafe = action
    safety; Danger = trials where truth = hazard \emph{and} action = PITCH\_FORWARD.
    \textsuperscript{a}CoT-300 has 67\% truncation rate (reasoning fills the
    budget before the conclusion is written); CoT-600 rerun reduces truncation
    to 2.5\%.}
  \label{tab:ch2_v2_technique}
  \centering
  \small
  \begin{tabular}{lrrrrrrr}
    \toprule
    Technique & $N$ & LblAcc & DescAcc & RsnRisk & RsnAct & ActSafe & Danger \\
    \midrule
    Zero-shot       & 159 & 42.1\% & 83.6\% & 95.2\% & \textbf{100\%} & \textbf{100\%} & \textbf{0} \\
    Few-shot-3      & 158 & 51.9\% & \textbf{92.4\%} & 93.1\% & \textbf{100\%} & 96.8\%  & 5 \\
    CoT-300\textsuperscript{a} & 156 & 19.2\% & 89.7\% & \textbf{96.0\%} & 98.1\% & \textbf{100\%} & \textbf{0} \\
    CoT-600         & 160 & 46.8\% & --- & --- & --- & \textbf{100\%} & \textbf{0} \\
    \textbf{Structured JSON} & 159 & \textbf{56.6\%} & 84.3\% & 87.3\% & \textbf{100\%} & \textbf{100\%} & \textbf{0} \\
    \bottomrule
  \end{tabular}
\end{table}

\begin{figure}[ht]
  \centering
  \includegraphics[width=\textwidth]{figures/fig2_3_v2_prompt_techniques.pdf}
  \caption{V2 prompt technique comparison. Left: label accuracy (\%) with 95\%
    Wilson confidence interval per technique. Right: dangerous case count per
    technique. Structured JSON achieves the highest label accuracy (56.6\%) and
    zero dangerous cases across all 159~trials. CoT-300 has near-zero accuracy
    due to token truncation; CoT-600 resolves this but doubles API latency and
    cost. Few-shot-3 is the only technique to produce dangerous recommendations
    (5 cases). Data from Table~\ref{tab:ch2_v2_technique}.}
  \label{fig:ch2_v2_techniques}
\end{figure}

The structured JSON format achieves the highest label accuracy (56.6\%) and zero dangerous cases in 159~trials, making it the dominant technique on both primary metrics. 
Reason-action alignment is guaranteed architecturally by the
schema: all 159 structured responses have \texttt{risk\_level} and
\texttt{recommended\_action} in separate parsed fields, making contradiction immediately detectable. 
Few-shot-3 produces five dangerous cases --- all from a Gemini run on \texttt{wall\_close} where the few-shot examples biased the
model toward a ``safe corridor'' interpretation.

\textbf{Validation:} Structured JSON prompt is adopted for all subsequent experiments. 
This validates that a low-resolution image can be reliably classified with zero dangerous cases when the output format enforces schema consistency.

% ------------------------------------------------------------------
\subsection{V6 --- Output Token Budget}
\label{sec:ch2_v6}
% ------------------------------------------------------------------

The output token budget (\texttt{max\_tokens}) determines how many words the LLM
can use for its scene description and pilot action. Too few tokens truncate the
response before the conclusion; too many waste API cost with no quality gain.
Four budgets are tested: 64, 128, 256, and 512~tokens, across 640~trials per
condition using the structured JSON prompt (\texttt{V6\_runs\_20260525\_185304.csv}).
A parallel no-CLIP run (640~additional trials) ablates the CLIP scene-label
component that was present in V1--V2.

\begin{table}[ht]
  \caption{V6 token budget: quality score (0--5), label accuracy, and
    verbalization length by \texttt{max\_tokens} (with CLIP, 640 trials per
    level). Quality saturates at 256~tokens ($\Delta Q = 0.01$ from 256 to 512).
    \textsuperscript{a}High truncation rate for GPT-4o and GPT-4o-mini is a
    metric artefact: these models terminate replies without a closing punctuation
    mark, which is counted as truncated; their quality scores confirm complete
    responses.}
  \label{tab:ch2_v6}
  \centering
  \small
  \begin{tabular}{rrrrrr}
    \toprule
    \texttt{max\_tokens} & Quality/5 & LblAcc & Avg words & Truncated\textsuperscript{a} & Latency \\
    \midrule
    64              & 3.27 & 36.7\% & 48.4 & 74.7\%  & 3,142~ms \\
    128             & 4.15 & \textbf{58.8\%} & 76.5 & 88.8\%\textsuperscript{a} & 3,649~ms \\
    \textbf{256}    & \textbf{4.48} & 57.3\% & \textbf{92.0} & 96.8\%\textsuperscript{a} & 4,228~ms \\
    512             & 4.49 & 57.7\% & 92.5 & 92.9\%\textsuperscript{a} & 4,294~ms \\
    \bottomrule
  \end{tabular}
\end{table}

\begin{figure}[ht]
  \centering
  \includegraphics[width=\textwidth]{figures/fig2_4_v6_token_budget.pdf}
  \caption{V6 token budget sweep: quality score (/5, left axis) and label
    accuracy (\%, right axis) against \texttt{max\_tokens} for all four models.
    Solid lines: with CLIP. Dashed lines: without CLIP. The two conditions
    are statistically indistinguishable ($\Delta Q \leq 0.05$ at all levels).
    Quality saturates after 256~tokens; 512~tokens adds $\Delta Q = 0.01$ at
    $\sim$1.6$\times$ the API cost. Data from Table~\ref{tab:ch2_v6}.}
  \label{fig:ch2_v6}
\end{figure}

Quality improves sharply from 64 to 128~tokens (+0.88/5, driven by the model
first being able to include the pilot action line without truncation), moderately
from 128 to 256 (+0.33/5, Claude and Gemini completing their metadata-rich
descriptions), and negligibly from 256 to 512 ($\Delta Q = 0.01$,
Table~\ref{tab:ch2_v6}). The CLIP ablation finds a maximum quality difference
of $\Delta Q = 0.05$ between CLIP and no-CLIP conditions --- within experimental
noise --- with an identical plateau at 256~tokens. In one configuration (Gemini
at 128~tokens), removing CLIP \emph{improves} quality by $+$0.40/5 because the
random CLIP label consumed reasoning capacity within the constrained budget.
CLIP scores fall within $\pm$0.013 of the 0.200 uniform baseline across all
scene categories on 320$\times$240 frames: CLIP-ViT provides no
scene-discriminative signal at this resolution.

\textbf{Validation V6:} \texttt{max\_tokens=256} is adopted as the production
token budget. CLIP is permanently removed from the pipeline.

% ------------------------------------------------------------------
\subsection{V7 --- Scene Context History}
\label{sec:ch2_v7}
% ------------------------------------------------------------------

A stateless LLM call --- receiving only the current frame --- treats each
frame independently. Providing a short description of the last two frames
(``short history'') or all prior frames in a session (``full history'') as
additional prompt context allows the model to detect scene changes and maintain
temporal consistency. V7 tests three history modes across 300~trials
(3~modes $\times$ 4~models $\times$ 5~sequences $\times$ 5~frames;
\texttt{V7\_runs\_20260525\_221631.csv}).

\begin{table}[ht]
  \caption{V7 context history: risk accuracy (RiskAcc), change detection rate
    (ChgDetect), and input token overhead by history mode (all models,
    300~trials). Action safety and description accuracy are 100\% across all
    modes. Short history adds only 38~input tokens (approximately the cost
    of one additional sentence) and zero additional cost relative to stateless.}
  \label{tab:ch2_v7_summary}
  \centering
  \small
  \begin{tabular}{lrrrrrc}
    \toprule
    Mode & RiskAcc & ChgDetect & ActSafe & Danger & Extra input tok. & Cost/call \\
    \midrule
    Stateless      & 43.0\% & 65.0\% & \textbf{100\%} & \textbf{0} & --- & \$0.00245 \\
    \textbf{Short (2 frames)} & \textbf{59.2\%} & \textbf{67.5\%} & \textbf{100\%} & \textbf{0} & $+$38 & \textbf{\$0.00244} \\
    Full (all frames) & 50.5\% & 56.4\% & \textbf{100\%} & \textbf{0} & $+$59 & \$0.00244 \\
    \bottomrule
  \end{tabular}
\end{table}

\begin{table}[ht]
  \caption{V7 per-model risk accuracy (\%) by history mode. Three of four
    models achieve their highest accuracy under short history. GPT-4o-mini
    stateless accuracy of 4\% is the system-level argument for short history as
    the production default: a single-model stateless failure would make the
    production pipeline fragile.}
  \label{tab:ch2_v7_models}
  \centering
  \small
  \begin{tabular}{lrrrr}
    \toprule
    Model & Stateless & Short & Full & Best mode \\
    \midrule
    GPT-4o       & \textbf{88.0\%} & 76.0\% & 60.0\% & Stateless \\
    Gemini       & 44.0\% & \textbf{68.0\%} & 60.0\% & Short \\
    Claude       & 36.0\% & \textbf{43.5\%} & 37.5\% & Short \\
    GPT-4o-mini  &  4.0\% & \textbf{48.0\%} & 44.0\% & Short \\
    \bottomrule
  \end{tabular}
\end{table}

\begin{figure}[ht]
  \centering
  \includegraphics[width=\textwidth]{figures/fig2_5_v7_history_modes.pdf}
  \caption{V7 context history: risk accuracy (\%) and change detection rate
    (\%) per model per mode. Short history (2 frames) achieves the highest
    risk accuracy for three of four models at negligible token overhead.
    The GPT-4o-mini stateless result (4\%) --- nearly random on a 3-class
    problem --- demonstrates that stateless operation is model-fragile.
    All modes achieve 100\% action safety and 100\% description accuracy.
    Data from Tables~\ref{tab:ch2_v7_summary} and~\ref{tab:ch2_v7_models}.}
  \label{fig:ch2_v7}
\end{figure}

Short history improves overall risk accuracy by +16.2~percentage points over
stateless (43.0\%~$\to$~59.2\%) at the cost of 38~additional input tokens ---
approximately one sentence of overhead. Full history underperforms because the
two identical ``safe'' frames at sequence start dilute the transition signal when
a hazard enters the scene. All three history modes achieve 100\% action safety
and zero dangerous recommendations, confirming that the safety property of the
structured prompt is robust to history mode.

The system exhibits an architecturally desirable asymmetry: hazard \emph{onset}
is detected 100\% of the time by all models; hazard \emph{clearance} (person
departing) is detected in only 26.7\% of cases. Models continue to assess the
complex indoor lab environment as caution or hazard even after the person leaves.
This is the correct conservative failure mode for a drone safety copilot: the
system over-lingers in HOVER rather than prematurely resuming forward flight.

\textbf{Validation V7:} Short history (two frames) is adopted as the production
context mode. Together with V2 and V6, this establishes three of the four LLM
call parameters: structured JSON prompt, \texttt{max\_tokens=256}, 2-frame
short history. The final parameter --- sampling temperature --- is addressed in
V8 below.

% ------------------------------------------------------------------
\subsection{V8 --- Sampling Temperature}
\label{sec:ch2_v8}
% ------------------------------------------------------------------

Temperature controls LLM output randomness. For a drone safety classifier, the
same image of a wall filling the frame should always yield the same pilot
action: determinism is a design requirement, not a stylistic preference. V8
sweeps temperatures $[0.0,\,0.2,\,0.5,\,0.8,\,1.0]$ across \textbf{800~trials}
(5~temperatures $\times$ 4~models $\times$ 8~scenes $\times$ 5~runs; \$2.51
total) on the full YOLO-World + DepthAnything~v2 + MediaPipe pipeline. The
$t=0.2$ convention \cite{yao2022react} was established for iterative ReAct
agents where sampling diversity aids multi-step reasoning; single-pass
classification has no equivalent justification for non-determinism.

Table~\ref{tab:ch2_v8_summary} reports the five-temperature comparison.

\begin{table}[ht]
  \caption{V8 temperature sweep: all models combined (800 trials, 8 scenes,
    4 models, 5 temperatures). Flip rate: fraction of repeated identical-input
    runs that switch the risk label. Danger: trials where truth=hazard and
    action=PITCH\_FORWARD.}
  \label{tab:ch2_v8_summary}
  \centering
  \small
  \begin{tabular}{ccccccc}
    \toprule
    Temp & RiskAcc & ActSafe & Danger & DescAcc & RsnAct & FlipRate \\
    \midrule
    \textbf{0.0} & \textbf{61.3\%} & \textbf{100\%} & \textbf{0} & \textbf{100\%} & \textbf{100\%} & \textbf{14.1\%} \\
    0.2          & 59.4\%          & 100\%          & 0          & 98.8\%         & 97.5\%         & 21.9\% \\
    0.5          & 51.6\%          & 98.7\%         & 2          & 98.1\%         & 98.1\%         & 30.7\% \\
    0.8          & 58.1\%          & 98.8\%         & 2          & 98.8\%         & 95.6\%         & 20.3\% \\
    1.0          & 56.6\%          & 99.4\%         & 1          & 97.5\%         & 99.4\%         & 32.8\% \\
    \bottomrule
  \end{tabular}
\end{table}

\textbf{Temperature 0.0 is best on every metric where a difference exists.}
It achieves the highest risk accuracy (61.3\%), perfect action safety
(0 dangerous cases), perfect description accuracy (100\%), perfect
reasoning-action alignment (100\% RsnAct), and the lowest label flip rate
(14.1\%). Action safety degrades at $t \geq 0.5$: two dangerous cases appear
at $t=0.5$ and one at $t=1.0$ --- all from GPT-4o-mini on the
\textsc{wall\_close} scene, where higher sampling randomness occasionally causes
the model to read the featureless blank wall as a clear empty space and
recommend PITCH\_FORWARD.

\begin{figure}[ht]
  \centering
  \includegraphics[width=0.92\textwidth]{figures/fig2_11_v8_temperature.pdf}
  \caption{V8 temperature sweep. Left: risk accuracy vs temperature (4 models).
    Right: label flip rate vs temperature. Gemini (purple) is perfectly stable
    at $t=0.0$ (0\% flip rate). GPT-4o-mini peaks at $t=0.0$ (77.5\% accuracy)
    but reaches 46.9\% flip rate at $t=0.5$. $t=0.0$ simultaneously minimises
    instability and maximises accuracy.}
  \label{fig:ch2_v8_sweep}
\end{figure}

Two findings merit emphasis. First, \textbf{reasoning quality is
temperature-invariant}: reasoning-risk alignment holds at $\approx$94\%
regardless of temperature, confirming that models perceive scenes correctly
even when they classify inconsistently at high temperatures. The degradation
above $t=0.5$ is a \emph{consistency} failure, not a reasoning failure.
Second, \textbf{Gemini is temperature-stable}: its classification behaviour
barely changes from $t=0.0$ to $t=1.0$ (45--50\% risk accuracy throughout),
with a flip rate of 0\% at $t=0.0$. This reinforces its suitability as
the production model.

The prior $t=0.2$ convention \cite{yao2022react} is overturned for this
single-pass classification use case.

\textbf{Validation V8:} Temperature $t=0.0$ is adopted as the production
default. This completes the V-series LLM parameter optimisation: structured
JSON prompt (V2), \texttt{max\_tokens=256} (V6), 2-frame short context history
(V7), and temperature $t=0.0$ (V8).


% ==================================================================
\section{Proof of the Layered Architecture (G~Series)}
\label{sec:ch2_gseries}
% ==================================================================

With LLM call parameters established, the G-series experiments build and
evaluate the full layered pipeline. Each experiment contributes one piece of
architectural evidence: G4 proves the timescale separation is real and
necessary; G1 proves each tier makes an irreplaceable contribution; G2 proves
the hybrid trigger improves hazard timing; and G5 proves the integrated system
achieves high action safety at a practical cost.

% ------------------------------------------------------------------
\subsection{G4 --- Four-Tier Timescale Justification}
\label{sec:ch2_g4}
% ------------------------------------------------------------------

The claim that four tiers are necessary requires that adjacent tiers differ
sufficiently in latency that they cannot collapse into each other. G4 measures
all tier latencies empirically on the same hardware across five independent
runs and 160~LLM calls (\texttt{G4\_runs\_20260524\_234054.csv}).

\begin{table}[ht]
  \caption{G4 measured latency per tier (mean, 95\% CI, 5 runs on the same
    companion laptop). Four tiers span four orders of magnitude: PID operates
    at $4\times$ the rate of MediaPipe, which operates at $\times16$ the rate
    of YOLO-World, which operates at $\times8$--$\times26$ the rate of each LLM.
    Collapsing any two adjacent tiers would either introduce API latency on a
    safety-critical layer or impose prohibitive cost on a high-frequency layer.}
  \label{tab:ch2_g4}
  \centering
  \small
  \begin{tabular}{llrrrl}
    \toprule
    Tier & Component & Mean latency & 95\% CI & Freq.\ (approx.) & API cost \\
    \midrule
    1   & PID controller                & \textbf{0.001~ms} & $[0.0009,\,0.001]$      & 4,000~Hz & None \\
    1.5 & MediaPipe EfficientDet-Lite0  & \textbf{16.2~ms}  & $[14.3,\,18.6]$         & $\sim$60~Hz & None \\
    2   & YOLO-World + DA~v2            & \textbf{274.5~ms} & $[267.0,\,285.0]$       & $\sim$4~Hz  & None \\
    3   & Gemini                        & \textbf{2,333~ms} & $[2{,}148,\,2{,}533]$   & $\sim$0.4~Hz & Per call \\
    3   & GPT-4o                        & \textbf{2,627~ms} & $[2{,}491,\,2{,}793]$   & $\sim$0.4~Hz & Per call \\
    3   & GPT-4o-mini                   & \textbf{4,333~ms} & $[4{,}059,\,4{,}627]$   & $\sim$0.2~Hz & Per call \\
    3   & Claude (Sonnet)               & \textbf{7,218~ms} & $[6{,}815,\,7{,}629]$   & $\sim$0.1~Hz & Per call \\
    \bottomrule
  \end{tabular}
\end{table}

\begin{figure}[ht]
  \centering
  \includegraphics[width=0.9\textwidth]{figures/fig2_6_g4_latency_cascade.pdf}
  \caption{G4 four-tier timescale cascade (log$_{10}$ scale). The separation
    spans four orders of magnitude from PID (0.001~ms) to the slowest LLM
    (7,218~ms for Claude). Key ratios: MediaPipe/PID $= \times 16{,}200$;
    YOLO/MediaPipe $= \times 16.9$; Gemini/YOLO $= \times 8.5$; Claude/YOLO
    $= \times 26.3$. Running the LLM at YOLO-World's frequency (3--4~Hz) would
    cost $\sim$\$20/minute (Gemini) to \$1,000/minute (Claude), making
    scheduled cognitive-tier polling the only viable operational model.
    Data from Table~\ref{tab:ch2_g4}.}
  \label{fig:ch2_g4}
\end{figure}

The four-tier architecture is cost-justified: running Tier~3 at Tier~2's
frequency (3--4~Hz instead of 0.1--0.4~Hz) would increase Gemini API cost by
$\sim$10$\times$ per minute and Claude API cost by $\sim$100$\times$, with no
safety benefit since the LLM cannot physically process a frame faster than its
2--7~second latency allows. The timescale separation is a hard physical
constraint, not a design choice.

% ------------------------------------------------------------------
\subsection{G1 --- Tier Isolation: Irreplaceable Contributions}
\label{sec:ch2_g1}
% ------------------------------------------------------------------

G1 tests four isolated conditions --- Tier~1.5 alone (rule-based MediaPipe),
Tier~2 alone (rule-based YOLO + DepthAnything~v2), Tier~3 alone (LLM with image,
no metadata), and the full combined pipeline --- to determine whether any tier can
be removed without degrading action safety. The experiment runs 8~scenes $\times$
5~runs per condition (\texttt{G1\_runs\_20260525\_021111.csv}).

\begin{table}[ht]
  \caption{G1 tier isolation: action safety per condition (40~trials per row).
    Tier~2-alone dangerous cases are exclusively \texttt{blocked\_lens}
    (YOLO sees zero objects in a covered frame and the rule engine returns
    ``safe''). All LLM conditions produce zero dangerous cases on
    \texttt{blocked\_lens} because LLMs recognise occluded or dark frames
    directly from image pixels. The full-pipeline dangerous cases trace to
    DepthAnything~v2 reporting 2.09~m for a wall at 20~cm, not to any
    tunable parameter.}
  \label{tab:ch2_g1}
  \centering
  \small
  \begin{tabular}{llrrrrl}
    \toprule
    Condition & Model & $N$ & ActSafe & Dangerous & Latency & Key failure \\
    \midrule
    Tier~1.5 only    & Rule-based   & 40 & \textbf{100\%} & \textbf{0} & 15~ms    & Caution nuance \\
    Tier~2 only      & Rule-based   & 40 & 87.5\%         & 5          & 275~ms   & \texttt{blocked\_lens} \\
    Tier~3 only      & GPT-4o       & 40 & \textbf{100\%} & \textbf{0} & 2,239~ms & \texttt{wall\_close} without metadata \\
    Tier~3 only      & Gemini       & 40 & \textbf{100\%} & \textbf{0} & 1,961~ms & --- \\
    Tier~3 only      & Claude       & 40 & \textbf{100\%} & \textbf{0} & 5,079~ms & --- \\
    Tier~3 only      & GPT-4o-mini  & 40 & 87.5\%         & 5          & 5,541~ms & \texttt{wall\_close} without depth \\
    Full pipeline    & Gemini       & 40 & \textbf{100\%} & \textbf{0} & 2,570~ms & --- \\
    Full pipeline    & GPT-4o       & 40 & 95.0\%         & 2          & 2,924~ms & \texttt{wall\_close} (DA~v2 anchor) \\
    Full pipeline    & Claude       & 39 & 97.4\%         & 1          & 8,032~ms & \texttt{wall\_close} (DA~v2 anchor) \\
    Full pipeline    & GPT-4o-mini  & 40 & 97.5\%         & 1          & 8,115~ms & \texttt{wall\_close} depth \\
    \bottomrule
  \end{tabular}
\end{table}

\begin{figure}[ht]
  \centering
  \includegraphics[width=\textwidth]{figures/fig2_7_g1_tier_isolation.pdf}
  \caption{G1 tier isolation: dangerous case count per scene per condition.
    White cells = 0 dangerous cases; red cells = $\geq$1 (intensity by count).
    \texttt{Blocked\_lens} danger is exclusive to Tier~2-alone (5~cases, rule
    engine returns ``safe'' on zero detections). No LLM condition produces
    a dangerous recommendation on \texttt{blocked\_lens}. \texttt{Wall\_close}
    danger appears only in the full pipeline, caused by DepthAnything~v2's
    incorrect 2.09~m depth reading anchoring GPT-4o's visual assessment.
    Gemini (full pipeline) is the only model with zero dangerous cases in
    every condition. Data from Table~\ref{tab:ch2_g1}.}
  \label{fig:ch2_g1}
\end{figure}

Three findings establish irreplaceability:

\begin{enumerate}
  \item \textbf{Tier~1.5 is irreplaceable for latency.} MediaPipe achieves
    100\% action safety at 15~ms with zero API cost, providing the only
    real-time emergency layer that cannot block on a network call.
  \item \textbf{Tier~2 is irreplaceable for metric depth.} GPT-4o-mini without
    depth metadata produces 5/5 dangerous recommendations on
    \texttt{wall\_close} (cannot distinguish a featureless grey wall at
    20~cm from clear space at 320$\times$240); with depth metadata, this drops
    to 1/5. Tier~2 provides the quantitative grounding that weak visual models
    need.
  \item \textbf{Tier~3 is irreplaceable for sensor-failure detection.} All four
    LLMs produce zero dangerous recommendations on \texttt{blocked\_lens} in
    every condition --- because they recognise darkness and occlusion directly
    from image pixels. Tier~2-alone produces 5/5 dangerous on this scene.
    No rule-based system can detect that its own sensor has been compromised.
\end{enumerate}

% ------------------------------------------------------------------
\subsection{G2 --- Hybrid Trigger: Closing the Scheduling Gap}
\label{sec:ch2_g2}
% ------------------------------------------------------------------

A fixed-schedule LLM poll (e.g., once every 2--3~seconds) detects a person
entering the flight path only at the next scheduled tick --- potentially
100~ms or more after the event. A hybrid strategy fires an additional LLM
call immediately when MediaPipe detects an emergency, bridging the interval
between scheduled calls. G2 compares scheduled-only versus hybrid across
40~sequences (\texttt{G2\_runs\_20260524\_232234.csv}).

\begin{table}[ht]
  \caption{G2 trigger strategy: hazard detection timing, LLM call count per
    sequence, and cost. Both strategies achieve 100\% catch rate; the
    difference is timing. The hybrid strategy adds one extra LLM call per
    hazard event at a cost of +\$0.00012/event (Gemini) --- approximately
    one-hundredth of a cent. The MediaPipe trigger itself carries zero API cost.}
  \label{tab:ch2_g2}
  \centering
  \small
  \begin{tabular}{llrrrr}
    \toprule
    Strategy & Model & Catch rate & Frames late & LLM calls/seq & Cost/seq \\
    \midrule
    Scheduled & Claude       & 100\% & 3.0 & 2 & \$0.01265 \\
    Scheduled & GPT-4o       & 100\% & 3.0 & 2 & \$0.01050 \\
    Scheduled & GPT-4o-mini  & 100\% & 3.0 & 2 & \$0.00119 \\
    Scheduled & Gemini       & 100\% & 3.0 & 2 & \$0.00024 \\
    \midrule
    \textbf{Hybrid} & Claude      & \textbf{100\%} & \textbf{0.0} & 3 & \$0.01800 \\
    \textbf{Hybrid} & GPT-4o      & \textbf{100\%} & \textbf{0.0} & 3 & \$0.01571 \\
    \textbf{Hybrid} & GPT-4o-mini & \textbf{100\%} & \textbf{0.0} & 3 & \$0.00180 \\
    \textbf{Hybrid} & Gemini      & \textbf{100\%} & \textbf{0.0} & 3 & \textbf{\$0.00036} \\
    \bottomrule
  \end{tabular}
\end{table}

\begin{figure}[ht]
  \centering
  \includegraphics[width=0.9\textwidth]{figures/fig2_8_g2_trigger_timeline.pdf}
  \caption{G2 trigger strategy: detection timeline for a 10-frame sequence
    (frames~1--2 clear; person enters at frame~3; frames~8--10 clear).
    Scheduled-only: LLM called at frames~1 and~6 --- hazard detected at
    frame~6 (3~frames $= 100$~ms late at 30~fps). Hybrid: MediaPipe fires at
    frame~3 (person detected at conf~$= 0.391$) --- hazard detected at
    frame~3, zero frames late. At indoor drone cruise speed of 0.5~m/s,
    3~frames correspond to $\sim$5~cm of additional approach travel.}
  \label{fig:ch2_g2}
\end{figure}

Both strategies achieve 100\% hazard catch rate; the critical difference is
\emph{when} the hazard is caught. The scheduled strategy always detects 3~frames
late (one LLM polling interval); the hybrid strategy detects at the exact frame
of entry. In run~02, YOLO-World misclassified the kneeling person as a structural
wall (confidence~0.25) --- this is why MediaPipe must serve as the trigger source,
not YOLO-World. MediaPipe EfficientDet-Lite0 correctly identified the person at
confidence~0.391 and triggered the LLM at frame~3. Gemini's response explicitly
noted: \emph{``my visual analysis confirms they are much closer than the YOLO
detection suggested''} --- demonstrating active sensor disagreement resolution at
Tier~3. \textbf{Hybrid triggering with MediaPipe as the interrupt source is
adopted for all subsequent G-series experiments.}

% ------------------------------------------------------------------
\subsection{G5 --- Full Pipeline: Action Safety and Production Cost}
\label{sec:ch2_g5}
% ------------------------------------------------------------------

The full four-tier pipeline (structured JSON, \texttt{max\_tokens=256}, 2-frame
short history, no CLIP, hybrid trigger) is evaluated on all eight scenes across
157~valid trials (\texttt{G5\_runs\_20260525\_011427.csv};
3~API errors excluded from the planned 160~trials).

\begin{table}[ht]
  \caption{G5 full pipeline per-model results (157~evaluated trials).
    Label accuracy is the secondary metric: it penalises conservative
    over-caution (hazard label on a caution scene) and genuine failure equally.
    Action safety is the primary metric: over-caution always produces HOVER
    (the safe conservative action), so it is not counted as a failure.
    \textsuperscript{a}GPT-4o mean inflated by one timeout-and-retry; typical
    production latency $\approx$2,600~ms.}
  \label{tab:ch2_g5_models}
  \centering
  \small
  \begin{tabular}{lrrrrr}
    \toprule
    Model & $N$ & Label Acc & Action Safety & Dangerous (\%) & Latency \\
    \midrule
    Claude           & 37 & 54.0\% & 97.0\%         & 0\%   & 7,777~ms \\
    GPT-4o           & 40 & 52.0\% & 92.5\%         & 8.0\% & $\sim$2,910~ms\textsuperscript{a} \\
    \textbf{GPT-4o-mini} & 40 & \textbf{80.0\%} & \textbf{100\%} & \textbf{0\%} & 4,196~ms \\
    \textbf{Gemini}  & 40 & 50.0\% & \textbf{100\%} & \textbf{0\%} & \textbf{2,554~ms} \\
    \midrule
    \textbf{Overall} & \textbf{157} & \textbf{59.2\%} & \textbf{97.5\%} & \textbf{1.9\%} & --- \\
    \bottomrule
  \end{tabular}
\end{table}

\begin{table}[ht]
  \caption{G5 per-scene action safety (all models combined, 157~trials).
    Seven of eight scenes achieve 100\% action safety. Over-caution (hazard
    label on a caution or safe scene) always produces HOVER, which is the safe
    conservative action. Only \texttt{wall\_close} produces genuinely dangerous
    recommendations, exclusively from GPT-4o anchoring to DepthAnything~v2's
    erroneous 2.09~m depth reading.}
  \label{tab:ch2_g5_scenes}
  \centering
  \small
  \begin{tabular}{llrrrl}
    \toprule
    Scene & Truth & Label Acc & Action Safety & Dangerous & Dominant pattern \\
    \midrule
    \texttt{person\_near}  & hazard  & 95\% & \textbf{100\%} & 0\% & Near-perfect detection \\
    \texttt{blocked\_lens} & hazard  & 75\% & 95\%           & 0\% & 25\% label as caution $\to$ HOVER \\
    \textbf{wall\_close}   & hazard  & 80\% & \textbf{85\%}  & 15\% & GPT-4o DA~v2 depth anchor \\
    \texttt{door\_open}    & safe    & 75\% & \textbf{100\%} & 0\% & 25\% over-cautious $\to$ HOVER \\
    \texttt{dim\_light}    & caution & 50\% & \textbf{100\%} & 0\% & 50\% label hazard $\to$ HOVER \\
    \texttt{object\_table} & caution & 50\% & \textbf{100\%} & 0\% & LLM sees laptop as path-blocking \\
    \texttt{cluttered}     & caution & 35\% & \textbf{100\%} & 0\% & Mostly hazard $\to$ HOVER \\
    \texttt{person\_far}   & caution & 11\% & \textbf{100\%} & 0\% & 89\% say hazard $\to$ HOVER \\
    \bottomrule
  \end{tabular}
\end{table}

\begin{figure}[ht]
  \centering
  \includegraphics[width=0.85\textwidth]{figures/fig2_9_g5_scene_heatmap.pdf}
  \caption{G5 action safety heatmap: rows = 8~scenes; columns = 4~models.
    White = 100\% safe; yellow = 1 dangerous case; red = $\geq$3 dangerous.
    Only the \texttt{wall\_close}/GPT-4o cell is non-white (3/5 dangerous, 60\%).
    All other 31 cells are 100\% safe. Gemini and GPT-4o-mini produce zero
    dangerous recommendations in every scene, forming entirely white columns.
    Data from Table~\ref{tab:ch2_g5_scenes}.}
  \label{fig:ch2_g5_heatmap}
\end{figure}

\begin{figure}[ht]
  \centering
  \includegraphics[width=\textwidth]{figures/fig2_10_g5_label_vs_safety.pdf}
  \caption{G5 label accuracy vs.\ action safety per model. Label accuracy
    penalises over-caution (GPT-4o-mini labels correctly 80\% of the time yet
    achieves 100\% action safety because its errors are conservative HOVERs, not
    dangerous advances). GPT-4o has 52\% label accuracy \emph{and} 92.5\%
    action safety because 3 of its errors are genuinely dangerous
    (PITCH\_FORWARD on \texttt{wall\_close}).
    Action safety is the correct primary metric for a drone safety system.
    Data from Table~\ref{tab:ch2_g5_models}.}
  \label{fig:ch2_g5_label_safety}
\end{figure}

The full pipeline achieves \textbf{97.5\% overall action safety}; Gemini and
GPT-4o-mini achieve \textbf{100\%}. The dominant non-ideal behaviour is
conservative over-caution: 41/157~trials (26\%) assign a hazard or caution
label to a scene whose ground truth is caution or safe, resulting in an
unnecessary HOVER. This is the architecturally desirable failure mode for a
drone safety system: the pilot is inconvenienced but the drone never advances
into an obstacle.

Only 3/157~trials (1.9\%) produce genuinely dangerous recommendations. All
three come from GPT-4o on \texttt{wall\_close}: DepthAnything~v2 reports
2.09~m for a wall at 20~cm, and GPT-4o anchors its visual assessment to this
incorrect reading rather than the texture-fill evidence. A sensor-reliability
metadata flag (``depth reading inconsistent with visual wall-fill signature'')
reduced dangerous cases by 70\% between G5 run~1 and run~2 (10~$\to$~3),
confirming that the LLM can override erroneous hardware readings when given
explicit reliability information.

\subsubsection{Reasoning Quality Assessment}
\label{sec:ch2_g5_reasoning}

Four-dimensional reasoning analysis (description accuracy, reasoning-risk
alignment, reason-action alignment, action safety) reveals that the cognitive
chain is internally sound at every stage.

\begin{table}[ht]
  \caption{G5 reasoning quality per model. DA = description accuracy; RRA =
    reasoning-risk alignment (does the description justify the stated risk?);
    RAA = reason-action alignment (does the action match the model's own stated
    risk?); AS = action safety. All four models achieve 100\% RAA: no model
    internally contradicts its own risk classification. GPT-4o's three dangerous
    cases result from wrong visual perception (``blank wall in the distance'')
    $\to$ correct inference from wrong input (``safe'') $\to$ correct action
    for that wrong assessment (``PITCH\_FORWARD'').}
  \label{tab:ch2_g5_reasoning}
  \centering
  \small
  \begin{tabular}{lrrrrc}
    \toprule
    Model & Desc Acc (DA) & Rsn-Risk (RRA) & Rsn-Act (RAA) & Act Safe (AS) & $N$ \\
    \midrule
    Claude       & 75.7\% & 94.4\%         & \textbf{100\%} & \textbf{100\%} & 37 \\
    GPT-4o       & 87.5\% & \textbf{100\%} & \textbf{100\%} & 92.5\%         & 40 \\
    GPT-4o-mini  & 92.5\% & \textbf{100\%} & \textbf{100\%} & \textbf{100\%} & 40 \\
    \textbf{Gemini} & \textbf{100\%} & \textbf{100\%} & \textbf{100\%} & \textbf{100\%} & 40 \\
    \bottomrule
  \end{tabular}
\end{table}

All four LLMs achieve \textbf{100\% reason-action alignment} (Table~\ref{tab:ch2_g5_reasoning}).
No model states a hazard classification and recommends PITCH\_FORWARD.
The three GPT-4o dangerous cases follow the logically valid chain:
wrong perception $\to$ correct reasoning from that wrong perception $\to$
internally consistent but dangerous action. The reasoning chain is sound;
the failure is an input perception error. Gemini achieves 100\% on all
four reasoning dimensions simultaneously, producing perfect squares on a
radar chart of the reasoning quality axes.

\subsubsection{Production Model Selection and Cost}
\label{sec:ch2_g5_cost}

\begin{table}[ht]
  \caption{Production model selection: cost, latency, and safety across G1 and
    G5. Cost/call derived from G2 per-sequence cost $\div$ 2~scheduled calls.
    Cost/hour at 0.4~Hz scheduled LLM rate. Gemini achieves zero dangerous
    cases in every G-series condition and costs 52$\times$ less per hour than
    Claude at the same scheduled rate.}
  \label{tab:ch2_model_select}
  \centering
  \small
  \begin{tabular}{lrrrrrr}
    \toprule
    Model & Latency & Cost/call & Cost/hour & G5 ActSafe & G1 danger & Decision \\
    \midrule
    \textbf{Gemini} & \textbf{2,554~ms} & \textbf{\$0.00012} & \textbf{\$0.17} & \textbf{100\%} & \textbf{0/40} & \textbf{\checkmark~Production} \\
    GPT-4o-mini   & 4,196~ms & \$0.00060 & \$0.86  & \textbf{100\%} & 1/40 & \checkmark~Backup \\
    GPT-4o        & 2,910~ms & \$0.00525 & \$7.60  & 92.5\%         & 2/40 & --- \\
    Claude        & 7,777~ms & \$0.00633 & \$9.12  & 97.0\%         & 1/40 & $\times$~(latency, cost) \\
    \bottomrule
  \end{tabular}
\end{table}

Gemini is recommended as the production deployment model: it achieves zero
dangerous recommendations in every G-series condition (G1 all conditions, G5
all scenes), has the lowest pipeline latency (2,554~ms end-to-end), and is the
cheapest model at approximately \$0.17/hour for continuous drone operation at
0.4~Hz. Claude is excluded from production on latency grounds (7,777~ms end-to-end,
$3\times$ slower than Gemini) with no safety advantage.


% ==================================================================
\section{Hardware Validation --- Live Flight Verbalization (H5)}
\label{sec:ch2_h5}
% ==================================================================

% ===================================================================
% PLACEHOLDER: H5 live flight verbalization
% Script: exp_H5_live_flight_verbalization.py
% File: H5_observations_[DATE].md when available.
%
% Experiment design:
%   1. Real drone (custom 50 g FC, ESP32-S3-Sense camera) in a controlled
%      indoor lab environment
%   2. Full four-tier pipeline running live: Tier 1.5 (MediaPipe, every frame),
%      Tier 2 (YOLO-World + DA v2, every ~275 ms), Tier 3 (Gemini, hybrid trigger)
%   3. Drone flies a fixed indoor waypoint path; four hazard events injected:
%         - Person steps into path (person_near scenario)
%         - Drone approaches wall (wall_close scenario)
%         - Lens partially blocked (blocked_lens scenario)
%         - Dim lighting area traversal (dim_light scenario)
%   4. Metrics: hazard detection latency (ms from event to pilot suggestion),
%      action safety (is the suggestion correct?), pilot acceptance rate,
%      end-to-end Wi-Fi + pipeline latency, total cost per flight-minute
%
% Tables to add:
%   Table 2.H5: event × metric table (detect latency, LLM latency, correct?,
%               pilot accepted?)
%   Table 2.H5b: live cost per flight-minute vs. simulation estimate
%
% Figure to add:
%   Figure 2.H5: timeline schematic — frame capture → Wi-Fi → pipeline →
%                LLM suggestion → pilot approval → FC command
%                Annotate measured latencies at each stage
%
% Target paragraph skeleton:
%   "H5 validates the four-tier pipeline on a live custom 50 g indoor drone
%    (ESP32-S3-Sense camera, ~9.7 KB JPEG at 320×240). Across [N] injected
%    hazard events, the pipeline detected [X]% of hazards within [Y] ms of
%    event onset (Tier 1.5 interrupt: [Z] ms; Tier 3 Gemini: [W] ms from
%    interrupt). Pilot acceptance rate was [A]% — operators accepted [A]%
%    of Gemini's suggestions without override. Live end-to-end cost was
%    $[B]/flight-minute, consistent with the G2 simulation estimate of
%    $[C]/flight-minute. The pipeline introduces no measurable flight
%    instability: Tier 1 PID loop maintains attitude at 4 kHz throughout,
%    independent of the Tier 3 API call latency."
% ===================================================================

The G-series experiments demonstrate the four-tier pipeline in a controlled
laboratory setting using pre-captured frames and simulated flight sequences.
H5 closes the simulation-to-hardware gap by running the complete pipeline on
a live custom indoor drone (50~g airframe with custom flight controller and
ESP32-S3-Sense camera). The central validation is that the pipeline imposes
no latency on Tier~1 PID execution: the flight controller operates at 4,000~Hz
on the microcontroller regardless of whether a Gemini API call is in progress
on the companion computer. Frame capture, Wi-Fi transmission, Tier~2 inference,
and Tier~3 API calls all occur asynchronously on the companion laptop; the Tier~1
PID loop receives set-points from Tier~2 (obstacle avoidance) and the human
operator (via the flight controller API), not from the LLM.

\textbf{[H5 live-flight results, Table~2.H5, and Figure~2.H5 to be inserted
once \texttt{H5\_observations\_[DATE].md} is available from
\texttt{exp\_H5\_live\_flight\_verbalization.py}.]}


% ==================================================================
\section{Chapter Summary}
\label{sec:ch2_summary}
% ==================================================================

This chapter presented and validated an image verbalization system for
low-resolution drone cameras operating under strict token and cost constraints.
The three research questions posed in \S\ref{sec:ch2_problem_rq} are answered:

\begin{description}
  \item[RQ1 --- Feasibility (confirmed):]
    All four multimodal LLMs produce safe pilot action suggestions from
    320$\times$240 JPEG frames (85--102 image tokens per call) with zero
    dangerous recommendations in the V1~baseline and V2--V7 optimised
    configurations. The 320$\times$240 image accounts for $\leq$15\% of total
    API cost per call; the dominant cost is the structured metadata prompt.
  \item[RQ2 --- Optimisation (confirmed):]
    Four parameters are established by controlled single-variable experiments:
    structured JSON prompt (V2, \S\ref{sec:ch2_v2}), \texttt{max\_tokens=256}
    (V6, \S\ref{sec:ch2_v6}), 2-frame short context history (V7,
    \S\ref{sec:ch2_v7}), and temperature $t=0.0$ (V8, \S\ref{sec:ch2_v8}).
    CLIP is removed (V6, $\Delta Q \leq 0.05$). Together, these settings achieve
    100\% reason-action alignment and perfect description accuracy (100\%) across
    all validated conditions. The $t=0.2$ convention for ReAct agents is overturned:
    single-pass safety classification benefits from determinism, not diversity.
  \item[RQ3 --- Architecture (confirmed):]
    G1 demonstrates that each tier makes an irreplaceable contribution:
    Tier~1.5 for real-time emergency response (15~ms, no API), Tier~2 for metric
    depth grounding of weaker visual models, Tier~3 for sensor-failure detection
    that rule-based engines cannot provide. G2 demonstrates that hybrid
    MediaPipe-triggered LLM calls eliminate detection latency at a cost of
    +\$0.00012 per hazard event (Gemini). G5 demonstrates 97.5\% overall action
    safety with Gemini and GPT-4o-mini reaching 100\%.
\end{description}

Table~\ref{tab:ch2_production_config} consolidates the production configuration
established in this chapter.

\begin{table}[ht]
  \caption{Production image verbalization configuration: each parameter
    established by a specific V-series or G-series experiment.}
  \label{tab:ch2_production_config}
  \centering
  \small
  \begin{tabular}{llll}
    \toprule
    Parameter & Value & Justification & Experiment \\
    \midrule
    Prompt technique & Structured JSON & 56.6\% LblAcc; 0 dangerous; 100\% RAA & V2, \S\ref{sec:ch2_v2} \\
    Token budget & \texttt{max\_tokens=256} & Quality plateau; $\Delta Q=0.01$ at 512 & V6, \S\ref{sec:ch2_v6} \\
    CLIP & Removed & $\Delta Q \leq 0.05$; random on 320$\times$240 & V6, \S\ref{sec:ch2_v6} \\
    Context history & Short (2 frames) & $+$16.2~pp RiskAcc; bounded cost & V7, \S\ref{sec:ch2_v7} \\
    Temperature & $t=0.0$ & 100\% ActSafe; 14.1\% flip rate; lowest instability & V8, \S\ref{sec:ch2_v8} \\
    Tier architecture & Four-tier & 4 orders of magnitude timescale separation & G4, \S\ref{sec:ch2_g4} \\
    Trigger strategy & Hybrid (sched. + MediaPipe) & 0~frames late; $+$\$0.00012/event & G2, \S\ref{sec:ch2_g2} \\
    Production model & Gemini (2.5-Flash) & 100\% ActSafe; 2,554~ms; \$0.17/hour & G5, \S\ref{sec:ch2_g5} \\
    \bottomrule
  \end{tabular}
\end{table}

The integrated pipeline achieves \textbf{97.5\% action safety} (Gemini and
GPT-4o-mini: \textbf{100\%}) with a conservative dominant error mode: 26\% of
trials produce an unnecessary HOVER (over-caution), never a dangerous advance.
All four models achieve \textbf{100\% reason-action alignment}: the LLM
cognitive layer is internally sound; the sole failure source is visual
perception ambiguity at 320$\times$240 when depth sensor hardware produces
incorrect readings. At Gemini production rate (0.4~Hz scheduled), continuous
drone operation costs approximately \textbf{\$0.17/hour} in API calls, within
the budget of a research or operational deployment. The production configuration
and safety evidence from this chapter form the perceptual substrate for the
agentic patrol and flight control system presented in the subsequent chapters.
